\label{sec:bias}

\subsection{The Maximum Minimum Representation Theorem}

The key property of Adams is that it maximises the minimum per-capita
representation across all states.

\begin{theorem}[Adams Maximises Minimum Representation]
  \label{thm:adams-minmax}
  Among all apportionments satisfying $\sum_i s_i = H$ and $s_i \geq 1$,
  Adams maximises $\min_i (p_i/s_i)$.
\end{theorem}

\begin{proof}
  For any apportionment $(s_1, \ldots, s_n)$, the minimum per-capita
  representation is $\min_i (p_i/s_i)$.
  Under Adams with divisor $d$, state $i$ receives $\lceil p_i/d \rceil$
  seats, so $s_i \leq p_i/d + 1$.
  Therefore $p_i/s_i \geq p_i/(p_i/d + 1) \geq d/(1 + d/p_i)$.

  For the state with the smallest population $p_{\min}$,
  $p_{\min}/s_{\min} \geq p_{\min} \cdot d / (p_{\min} + d) \to d$
  as $d \to 0$.
  The Adams divisor is chosen to achieve exactly $H$ total seats;
  the minimum representation is determined by the smallest population
  state's ceiling.

  Alternative proof: Suppose some apportionment $(s_1', \ldots, s_n')$
  with the same total $H$ achieves higher minimum representation
  $\min_i (p_i/s_i') > \min_i (p_i/s_i^{\text{Ad}})$.
  Let $j = \arg\min_i (p_i/s_i^{\text{Ad}})$ be the state with lowest
  representation under Adams.
  Since $p_j/s_j' > p_j/s_j^{\text{Ad}}$, we have $s_j' < s_j^{\text{Ad}}$.
  But $s_j^{\text{Ad}} = \lceil p_j/d \rceil$ is already the smallest
  possible value for a ceiling-rounding method at divisor $d$.
  Any smaller $s_j'$ would violate ceiling rounding for some achievable
  divisor, contradicting the assumption.
  See \citet{balinski1982fair}, ch.~3, for a rigorous treatment.
\end{proof}

\subsection{Equivalent Formulation: Upper Quota Guarantee}

Adams guarantees that every state receives at least $\lceil q_i \rceil$
seats (upper quota), where $q_i = p_i \times H / N$ is the exact
proportional quota.
This is equivalent to Theorem~\ref{thm:adams-minmax}:
if every state receives its upper quota, then the most underrepresented
state (the one most harmed by fractional rounding) is maximally protected.

\begin{corollary}[Upper Quota Compliance]
  Under Adams, every state satisfies $s_i \geq \lceil q_i \rceil$,
  i.e., no state receives fewer seats than its exact proportional
  entitlement rounded up.
\end{corollary}

This is a strong form of fairness from the perspective of small states:
no state is ``cheated'' of any fractional entitlement.
The cost is that large states are correspondingly underrepresented
relative to their exact quotas.

\subsection{The Bias Ordering}

For a state with exact quota $q_i = 1 + \epsilon$ (barely above 1):
\begin{itemize}
  \item Adams: receives 2 seats (ceiling rounding).
  \item Huntington-Hill: receives 2 seats if $\epsilon > \sqrt{2} - 1
    \approx 0.414$; 1 seat otherwise.
  \item Webster: receives 2 seats if $\epsilon > 0.5$; 1 seat otherwise.
  \item Jefferson: receives 1 seat (floor rounding, regardless of $\epsilon$).
\end{itemize}
This comparison makes the bias ordering concrete:
Adams gives a bonus seat to states with \emph{any} fractional quota
above 1; Jefferson gives no bonus seat regardless of the fractional part;
HH and Webster occupy intermediate positions.
For small states with $q_i$ near 1 (the most common case where methods
disagree), Adams is maximally generous and Jefferson is maximally
stingy.
